% ID: FIS-DIN-NEW-003
% Titolo: Tempo e velocita lungo la direzione della risultante
% Disciplina: Fisica
% Area: Dinamica
% Argomento: Leggi di Newton
% Sottoargomento: Accelerazione lungo la risultante delle forze
% Classe: Terza liceo scientifico
% Difficolta: 4
% Tipo: problema_numerico
% Risultato: F_tot = (12,2; -3,0) N; t = 1,69 s; v = 17,7 m/s
% Tempo_stimato: 12 min
% Competenze: scomposizione vettoriale, dinamica, cinematica lungo una direzione
% Prerequisiti: secondo principio della dinamica, vettori, moto uniformemente accelerato
% Tag: fisica, dinamica, leggi_newton, risultante, moto_accelerato
% Fonte: verifica_fisica_dinamica_anonimizzata
% Autore: Riccardo Carli
% Licenza: CC-BY-SA-4.0

\begin{esercizio}
Un corpo di massa \(m = 1{,}2\,\mathrm{kg}\) parte da fermo ed e soggetto a due
forze costanti: \(F_1 = 7{,}0\,\mathrm{N}\), diretta lungo \(+\hat x\), e
\(F_2 = 6{,}0\,\mathrm{N}\), diretta a \(-30^\circ\) rispetto allo stesso
asse.

Dopo aver calcolato la forza totale
\[
\vec F_\text{tot} =
\begin{pmatrix}
F_x \\[4pt]
F_y
\end{pmatrix},
\]
e la corrispondente accelerazione, determina il tempo necessario a percorrere
una distanza di \(15\,\mathrm{m}\) lungo la direzione della risultante. Indica
infine la velocita del corpo al termine del tragitto.
\end{esercizio}

\begin{soluzione}
La prima forza ha componenti:
\[
\vec F_1=
\begin{pmatrix}
7{,}0\\[4pt]
0
\end{pmatrix}
\mathrm{N}.
\]

La seconda forza e diretta a \(-30^\circ\), quindi:
\[
F_{2x}=6{,}0\cos(-30^\circ)\simeq 5{,}20\,\mathrm{N},
\]
\[
F_{2y}=6{,}0\sin(-30^\circ)=-3{,}0\,\mathrm{N}.
\]

La forza totale e:
\[
\vec F_\text{tot}=
\begin{pmatrix}
7{,}0+5{,}20\\[4pt]
0-3{,}0
\end{pmatrix}
=
\begin{pmatrix}
12{,}2\\[4pt]
-3{,}0
\end{pmatrix}
\mathrm{N}.
\]

Il modulo della risultante e:
\[
|\vec F_\text{tot}|
=\sqrt{(12{,}2)^2+(-3{,}0)^2}
\simeq 12{,}6\,\mathrm{N}.
\]

L'accelerazione vettoriale e:
\[
\vec a=\frac{\vec F_\text{tot}}{m}
=
\begin{pmatrix}
10{,}2\\[4pt]
-2{,}5
\end{pmatrix}
\mathrm{m/s^2}.
\]
Il suo modulo e:
\[
a=\frac{|\vec F_\text{tot}|}{m}
=\frac{12{,}6}{1{,}2}
\simeq 10{,}5\,\mathrm{m/s^2}.
\]

Poiche il corpo parte da fermo, il moto avviene lungo la direzione della
risultante. Lungo tale direzione:
\[
s=\frac{1}{2}at^2.
\]
Con \(s=15\,\mathrm{m}\):
\[
t=\sqrt{\frac{2s}{a}}
=\sqrt{\frac{30}{10{,}5}}
\simeq 1{,}69\,\mathrm{s}.
\]

La velocita finale lungo la direzione della risultante e:
\[
v=at=10{,}5\cdot 1{,}69
\simeq 17{,}7\,\mathrm{m/s}.
\]
\end{soluzione}
